\subsection{LLM valutati}

Il sistema sviluppato in questo lavoro è stato valutato con quattro LLM open-weight, scelti tra
quelli eseguibili sull'hardware a disposizione. Per tre di essi questo è possibile solo ricorrendo a una
quantizzazione. La \emph{quantizzazione} riduce la precisione numerica con cui sono memorizzati i pesi di un modello,
per esempio da 16 a 4 bit per peso: la memoria occupata si riduce di circa quattro volte, al prezzo di una possibile
perdita di qualità delle risposte. I tre modelli quantizzati sono serviti nella variante \emph{AWQ}
(\emph{Activation-aware Weight Quantization}), uno
schema di quantizzazione post-addestramento a 4 bit proposto da Lin et al.\ \cite{lin2024awq}. Gli autori osservano
che una piccola frazione dei pesi di un modello, quelli associati ai canali di attivazione di ampiezza maggiore,
contribuisce in modo sproporzionato all'errore di quantizzazione. AWQ individua questi pesi salienti a partire dalle
statistiche di attivazione, anziché dalla sola magnitudine dei pesi, e li protegge applicando un fattore di scala
per canale prima della quantizzazione: in questo modo la perdita di qualità si riduce drasticamente rispetto a uno
schema uniforme, senza richiedere un nuovo addestramento né una fase di backpropagation. Il quarto modello,
Llama-3.1-8B, è invece eseguito nella versione originale non quantizzata, con pesi in formato BF16 (numeri a 16 bit, la precisione con cui il modello è distribuito).

I quattro modelli confrontati sono:

\begin{itemize}
  \item \textbf{Qwen2.5-14B-Instruct-AWQ} e \textbf{Qwen2.5-7B-Instruct-AWQ}, due taglie della famiglia Qwen2.5
    sviluppata da Alibaba Cloud \cite{qwen2024technical}, che differiscono di un fattore 2 nel numero di parametri
    (14 contro 7 miliardi). Gli autori riportano un corpus di pre-training di 18 mila miliardi di token, più del
    doppio dei 7 mila miliardi della generazione precedente (Qwen2), con particolare attenzione a codice e
    matematica, e un supporto multilingue esteso a oltre 29 lingue, italiano compreso. Entrambe le varianti posseggono una finestra di contesto
    di 128k token;
  \item \textbf{Mistral-Nemo-12B-AWQ}, sviluppato congiuntamente da Mistral AI e NVIDIA
    \cite{mistralai2024nemo}, un modello da 12 miliardi di parametri con una finestra di contesto di 128k token.
    Introduce un nuovo tokenizer, Tekken, che secondo gli autori rappresenta il testo con circa il 30\% di token in
    meno rispetto al tokenizer dei precedenti modelli Mistral su codice sorgente e su diverse lingue europee, tra
    cui l'italiano, e risulta più efficiente del tokenizer di Llama~3 per circa l'85\% delle lingue. Il modello è
    stato inoltre addestrato con \emph{quantization awareness}, una procedura pensata per consentire l'inferenza in
    FP8, cioè con pesi a 8 bit;
  \item \textbf{Meta-Llama-3.1-8B-Instruct}, appartenente alla famiglia Llama~3 di Meta \cite{grattafiori2024llama3},
    composta da tre taglie dense (8, 70 e 405 miliardi di parametri) addestrate su oltre 15 mila miliardi di token,
    con una finestra di contesto di 128k token. Gli autori dichiarano un supporto ufficiale per otto lingue,
    italiano compreso, e un addestramento specifico per l'uso di strumenti esterni (\emph{tool calling}), una
    capacità agentica non rilevante per l'uso puramente conversazionale che ne fa questo sistema.
\end{itemize}

Tutti e quattro i modelli condividono la stessa architettura di fondo, il Transformer decoder-only auto-regressivo, e differiscono principalmente per numero di parametri, dati e
procedura di addestramento e tokenizer, oltre che per il formato dei pesi: Llama-3.1-8B è l'unico eseguito senza
quantizzazione, e di questa differenza occorre tenere conto nell'interpretare il confronto, perché una parte del
divario di qualità o di latenza rispetto agli altri modelli può dipendere dal formato più che dal modello in sé.
La Tabella~\ref{tab:llm_comparison} ne riassume le caratteristiche principali, mentre il confronto sistematico in termini di qualità della risposta generata, latenza e comportamento sulle diverse
categorie del dataset di test è presentato nel Capitolo~\ref{cap:validazione}.

\begin{table}[H]
    \centering
    \footnotesize
    \setlength{\tabcolsep}{4pt}
    \renewcommand{\arraystretch}{1.35}
    \begin{tabularx}{\textwidth}{@{} L{3.4cm} C{1.2cm} C{1.6cm} C{1.4cm} C{1.4cm} >{\raggedright\arraybackslash}X @{}}
        \toprule
        \textbf{Modello} & \textbf{Param.} & \textbf{Contesto} & \textbf{Formato} & \textbf{VRAM (pesi)} & \textbf{Caratteristiche principali} \\
        \midrule
        \textbf{Qwen2.5-14B-Instruct-AWQ}
            & 14,7 mld
            & 128k
            & AWQ 4 bit
            & $\sim$9--10 GB
            & Pre-training su 18 mila miliardi di token; oltre 29 lingue, italiano compreso. \\
        \midrule
        \textbf{Qwen2.5-7B-Instruct-AWQ}
            & 7,6 mld
            & 128k
            & AWQ 4 bit
            & $\sim$5--6 GB
            & Stessa famiglia del 14B con circa metà dei parametri: minore latenza e occupazione di memoria. \\
        \midrule
        \textbf{Mistral-Nemo-12B-AWQ}
            & 12,2 mld
            & 128k
            & AWQ 4 bit
            & $\sim$8--9 GB
            & Tokenizer Tekken, più compatto sull'italiano rispetto ai precedenti tokenizer Mistral. \\
        \midrule
        \textbf{Meta-Llama-3.1-8B-Instruct}
            & 8,0 mld
            & 128k
            & BF16 \newline
            & $\sim$16 GB
            & Pre-training su oltre 15 mila miliardi di token; otto lingue supportate, italiano compreso. \\
        \bottomrule
    \end{tabularx}
    \caption{Caratteristiche dei quattro modelli generativi valutati.}
    \label{tab:llm_comparison}
\end{table}